Since the computation is large, we made a few sanity checks of the program and of its output;
the scripts are in the repository, with the database. For $n\le10$, an exact test of every
connected graph on $n$ vertices from nauty's \ml{geng -c} finds exactly the graphs of $\cL_n$. A
second program rereads the files and checks, with PARI~\cite{PARI}, that every listed graph is
exceptional, in the right list and listed once. A SageMath script finds the seven exceptional
strongly regular graphs~\cite[Theorem~3.12.4]{BCN89} in list (A), and none of $34$ small
generalized line graphs in the lists.

\subsection*{Comparison with the literature}
List (B) has, on $6$, $7$ and $8$ vertices, the numbers of graphs that~\cite[\S2.3 and
Table~A2]{CRS04} and~\cite{DC79} give for the exceptional graphs with $\lmin>-2$
(Section~\ref{sec:intro}), and the numbers of maximal exceptional graphs (Section~\ref{sec:counts}) agree
with~\cite{CLRS02} and~\cite{Cve04}. Cvetkovi\'c
continues the passage of~\cite{Cve04} quoted in Section~\ref{sec:intro} with ``Minimal
exceptional graphs, all having 6 vertices, were found some 25 years ago'' (see
also~\cite{CRS04}). Every exceptional graph on $6$ vertices is minimal, since an exceptional
graph has at least $6$ vertices by Proposition~\ref{prop:types}; thus the minimal exceptional
graphs are exactly the exceptional graphs on $6$ vertices, and there are $20$ of them in list
(B), as~\cite{Cve04} states.

Cvetkovi\'c also writes that ``Most of the maximal exceptional graphs (432 of them) are 29 vertex
cones''~\cite{Cve04}. The maximal exceptional graphs are divided in~\cite[Chapter~6]{CRS04}, as
reported in~\cite[\S3]{DBvV26}, into three types: (a) those with largest valency $28$, the
largest valency possible~\cite[Theorem~4.3]{CGSS}, and exactly $29$ vertices; (b) those with
largest valency $28$ and more than $29$ vertices; and (c) those with largest valency less than
$28$. De Bruyn and van Veluw state that there are $430$ graphs of type (a), $37$ of type (b) and
$6$ of type (c); a graph of type (a) is a cone (a graph with a vertex adjacent to all the others),
since its vertex of valency $28$ is adjacent to all the others. Our counts agree exactly: $430$ of
our maximal exceptional graphs have $29$ vertices and a vertex of valency $28$, $37$ have more
than $29$ vertices and a vertex of valency $28$, and $6$ have largest valency less than $28$. Two
of the last $6$ have $29$ vertices, and both have largest valency $26$:
\begin{center}
\small
\verb!\??GWWJ@}[ThijhfxLuIzom}`vx{]AyXdYijTeUxrLNEiizKkrnS[N^}@~~{~~~~Z~~~{!\\
\verb!\@G?GGBEDkLp\R[jthtTZkY|hfyk}C[xeTYjcmUyYlNHUizS[rna[N^~B~~}N^~~N~~~{!
\end{center}
Thus exactly $430$ of the $432$ maximal exceptional graphs on $29$ vertices are cones: the
number $432$ of~\cite{Cve04} is the number of all maximal exceptional graphs on $29$ vertices.

The program \ml{count_regular.py} finds $187$ regular graphs in the lists, all in list (A), in
agreement with the $187$ regular exceptional graphs of Bussemaker, Cvetkovi\'c and
Seidel~\cite{BCS} (see~\cite[\S3.11]{BCN89} and~\cite{DBvV26}). With $k$ the valency, $163$ of
them have $n=2(k+2)$, $21$ have $n=3(k+2)/2$ and $3$ have $n=4(k+2)/3$, as they must
by~\cite[Theorem~3.12.2]{BCN89}. Hoffman's ratio bound on the independence number of a
$k$-regular graph on $n$ vertices with least eigenvalue $-2$ is $2n/(k+2)$, so these are the
classes of $163$, $21$ and $3$ regular exceptional graphs of~\cite{DBvV26}, with ratio bounds
$4$, $3$ and $8/3$.
